% GENERATED FILE. Edit canonical lessons, then run npm run generate:books.
% canonical-source-hash: fnv1a64:2c5313e99be0b3bd
% canonical-lessons: MR-C166-vaya-ghalavne, MR-C166-gappa-basne, MR-C166-ghabarne, MR-C166-vahun-nene, MR-C166-marne

\chapter{To waste, To keep quiet, To be afraid, To carry, To hit}
\label{ch:mr-to-waste-to-keep-quiet-to-be-afraid-to-carry-to-hit}
\hlchaptermodality{166}

\begin{chapteropening}
\textbf{By the end of this chapter:} \emph{I can say to waste, to keep quiet, to be afraid, to carry and to hit.}

Complete the last lesson of chapter 166: I can say to waste, to keep quiet, to be afraid, to carry and to hit.
\end{chapteropening}

\section[\texorpdfstring{\mr{वाया} \mr{घालवणे} (vāyā ghālavṇe)}{वाया घालवणे (vāyā ghālavṇe)}]{\mr{वाया} \mr{घालवणे} (vāyā ghālavṇe) — to waste}

\label{lesson:MR-C166-vaya-ghalavne}



\begin{quote}
Before the new one: say the Marathi for to observe, then the Marathi for to take part.
\end{quote}

\subsection*{You'll want to know: \mr{वाया} \mr{घालवणे}}
\textbf{\mr{वाया} \mr{घालवणे}} — \emph{vāyā ghālavṇe} — \textquotedblleft{}to waste\textquotedblright{}.

\begin{cousinweb}
The origin of \textbf{\mr{वाया}}, \textquotedblleft{}in vain\textquotedblright{}, is not settled: CDIAL has no entry for it. \textbf{\mr{घालवणे}} is the \textquotedblleft{}make\textquotedblright{} form of \textbf{\mr{घालणे}}, which CDIAL lists, as \textquotedblleft{}to pour\textquotedblright{}, under a reconstructed \emph{*ghalyati}, \textquotedblleft{}pours\textquotedblright{} (CDIAL 4457).
\end{cousinweb}

\begin{grammarlens}[title={What you've built}]
One more word for this chapter.
\end{grammarlens}

\subsection*{Your turn}
\begin{itemize}
\raggedright
  \item \emph{Say it:} \emph{vāyā ghālavṇe}
  \item \emph{Say it:} \emph{vāyā ghālavṇe}, once more
  \item \emph{Say it:} say \emph{vāyā ghālavṇe}
  \item \emph{Recall:} say the Marathi for to observe, then the Marathi for to take part, then say \emph{vāyā ghālavṇe} again
\end{itemize}

\begin{tcolorbox}[breakable,colback=teal!4,colframe=teal!35!black,title={Before you move on}]
What does \mr{वाया} \mr{घालवणे} mean? (\textquotedblleft{}To waste\textquotedblright{}.) Say it once more.
\end{tcolorbox}
\section[\texorpdfstring{\mr{गप्प} \mr{बसणे} (gappa basṇe)}{गप्प बसणे (gappa basṇe)}]{\mr{गप्प} \mr{बसणे} (gappa basṇe) — to keep quiet}

\label{lesson:MR-C166-gappa-basne}



\begin{quote}
Before the new one: say the Marathi for to take part, then the Marathi for to waste.
\end{quote}

\subsection*{You'll want to know: \mr{गप्प} \mr{बसणे}}
\textbf{\mr{गप्प} \mr{बसणे}} — \emph{gappa basṇe} — \textquotedblleft{}to keep quiet\textquotedblright{}.

\begin{cousinweb}
The origin of \textbf{\mr{गप्प}}, \textquotedblleft{}silent\textquotedblright{}, is not settled: CDIAL has no entry for it in this sense. \textbf{\mr{बसणे}} comes from Sanskrit \emph{upaviśati}, \textquotedblleft{}sits down\textquotedblright{} (CDIAL 2245).
\end{cousinweb}

\begin{grammarlens}[title={What you've built}]
One more word for this chapter.
\end{grammarlens}

\subsection*{Your turn}
\begin{itemize}
\raggedright
  \item \emph{Say it:} \emph{gappa basṇe}
  \item \emph{Say it:} \emph{gappa basṇe}, once more
  \item \emph{Say it:} say \emph{gappa basṇe}
  \item \emph{Recall:} say the Marathi for to take part, then the Marathi for to waste, then say \emph{gappa basṇe} again
\end{itemize}

\begin{tcolorbox}[breakable,colback=teal!4,colframe=teal!35!black,title={Before you move on}]
What does \mr{गप्प} \mr{बसणे} mean? (\textquotedblleft{}To keep quiet\textquotedblright{}.) Say it once more.
\end{tcolorbox}
\section[\texorpdfstring{\mr{घाबरणे} (ghābarṇe)}{घाबरणे (ghābarṇe)}]{\mr{घाबरणे} (ghābarṇe) — to be afraid}

\label{lesson:MR-C166-ghabarne}



\begin{quote}
Before the new one: say the Marathi for to waste, then the Marathi for to keep quiet.
\end{quote}

\subsection*{You'll want to know: \mr{घाबरणे}}
\textbf{\mr{घाबरणे}} — \emph{ghābarṇe} — \textquotedblleft{}to be afraid\textquotedblright{}.

\begin{cousinweb}
The origin of \textbf{\mr{घाबरणे}} (\emph{ghābarṇe}) is not settled. Neither CDIAL nor DEDR, searched by form and by meaning, has an entry for it.
\end{cousinweb}

\begin{grammarlens}[title={What you've built}]
One more word for this chapter.
\end{grammarlens}

\subsection*{Your turn}
\begin{itemize}
\raggedright
  \item \emph{Say it:} \emph{ghābarṇe}
  \item \emph{Say it:} \emph{ghābarṇe}, once more
  \item \emph{Say it:} say \emph{ghābarṇe}
  \item \emph{Recall:} say the Marathi for to waste, then the Marathi for to keep quiet, then say \emph{ghābarṇe} again
\end{itemize}

\begin{tcolorbox}[breakable,colback=teal!4,colframe=teal!35!black,title={Before you move on}]
What does \mr{घाबरणे} mean? (\textquotedblleft{}To be afraid\textquotedblright{}.) Say it once more.
\end{tcolorbox}
\section[\texorpdfstring{\mr{वाहून} \mr{नेणे} (vāhūn neṇe)}{वाहून नेणे (vāhūn neṇe)}]{\mr{वाहून} \mr{नेणे} (vāhūn neṇe) — to carry, to transport}

\label{lesson:MR-C166-vahun-nene}



\begin{quote}
Before the new one: say the Marathi for to keep quiet, then the Marathi for to be afraid.
\end{quote}

\subsection*{You'll want to know: \mr{वाहून} \mr{नेणे}}
\textbf{\mr{वाहून} \mr{नेणे}} — \emph{vāhūn neṇe} — \textquotedblleft{}to carry, to transport\textquotedblright{}.

\begin{cousinweb}
\textbf{\mr{वाहून}} is \textquotedblleft{}having carried\textquotedblright{}, from Marathi \emph{vāhṇẽ}, \textquotedblleft{}to carry\textquotedblright{}, which CDIAL lists under Sanskrit \emph{vāhayati}, \textquotedblleft{}drives (a chariot)\textquotedblright{} (CDIAL 11612). \textbf{\mr{नेणे}} comes from Sanskrit \emph{nayati}, \textquotedblleft{}leads, carries off\textquotedblright{} (CDIAL 6966).
\end{cousinweb}

\begin{grammarlens}[title={What you've built}]
One more word for this chapter.
\end{grammarlens}

\subsection*{Your turn}
\begin{itemize}
\raggedright
  \item \emph{Say it:} \emph{vāhūn neṇe}
  \item \emph{Say it:} \emph{vāhūn neṇe}, once more
  \item \emph{Say it:} say \emph{vāhūn neṇe}
  \item \emph{Recall:} say the Marathi for to keep quiet, then the Marathi for to be afraid, then say \emph{vāhūn neṇe} again
\end{itemize}

\begin{tcolorbox}[breakable,colback=teal!4,colframe=teal!35!black,title={Before you move on}]
What does \mr{वाहून} \mr{नेणे} mean? (\textquotedblleft{}To carry, to transport\textquotedblright{}.) Say it once more.
\end{tcolorbox}
\section[\texorpdfstring{\mr{मारणे} (mārṇe)}{मारणे (mārṇe)}]{\mr{मारणे} (mārṇe) — to hit; to kill}

\label{lesson:MR-C166-marne}



\begin{quote}
Before the new one: say the Marathi for to be afraid, then the Marathi for to carry.
\end{quote}

\subsection*{You'll want to know: \mr{मारणे}}
\textbf{\mr{मारणे}} — \emph{mārṇe} — \textquotedblleft{}to hit; to kill\textquotedblright{}.

\begin{cousinweb}
\textbf{\mr{मारणे}} (\emph{mārṇe}) is a tadbhava, a word inherited from Sanskrit and reshaped by sound change. It comes from \emph{mārayati}, \textquotedblleft{}kills\textquotedblright{} (CDIAL 10066). It is the \textquotedblleft{}make\textquotedblright{} form of \emph{mṛ}, \textquotedblleft{}to die\textquotedblright{}: to kill is to make die. CDIAL glosses the Marathi verb \textquotedblleft{}to strike, kill\textquotedblright{}.
\end{cousinweb}

\begin{grammarlens}[title={What you've built}]
One more word for this chapter.
\end{grammarlens}

\subsection*{Your turn}
\begin{itemize}
\raggedright
  \item \emph{Say it:} \emph{mārṇe}
  \item \emph{Say it:} \emph{mārṇe}, once more
  \item \emph{Say it:} say \emph{mārṇe}
  \item \emph{Recall:} say the Marathi for to be afraid, then the Marathi for to carry, then say \emph{mārṇe} again
\end{itemize}

\begin{tcolorbox}[breakable,colback=teal!4,colframe=teal!35!black,title={Before you move on}]
What does \mr{मारणे} mean? (\textquotedblleft{}To hit; to kill\textquotedblright{}.) Say it once more.
\end{tcolorbox}
